% ============================================
% 山西大学第二十届数学建模竞赛论文模板（校赛版）
%
% 校赛要求(依据附件1)：
%   1. 首页：论文题目 + 摘要（含模型主要特点、建模方法、主要结果）+ 关键词
%   2. 正文：问题提出 -> 模型假设 -> 模型建立 -> 模型求解 -> 结果分析 -> 模型展望
%   3. 参考文献(标准格式)
%   4. 附录：计算机程序
% 提交：PDF格式，命名"题号+姓名"，如 A_张三_李四_王五
% 编译方式: xelatex -> xelatex (2次)
% ============================================
\documentclass[12pt,a4paper]{article}

% ---------- 中文支持 ----------
\usepackage{ctex}
\usepackage{fontspec}
\setmainfont{Times New Roman}

% ---------- 数学公式 ----------
\usepackage{amsmath,amssymb,amsfonts}
\usepackage{bm}

% ---------- 页面设置 ----------
\usepackage{geometry}
\geometry{left=2.5cm,right=2.5cm,top=2.5cm,bottom=2.5cm}

% ---------- 图表 ----------
\usepackage{graphicx}
\usepackage{booktabs}
\usepackage{longtable}
\usepackage{float}
\usepackage{subfigure}

% ---------- 超链接 ----------
\usepackage{hyperref}
\hypersetup{colorlinks=true, linkcolor=black, citecolor=black, urlcolor=blue}

% ---------- 页眉页脚 ----------
\usepackage{fancyhdr}
\pagestyle{fancy}
\setlength{\headheight}{14.5pt}
\lhead{山西大学第二十届数学建模竞赛}
\chead{}
\rhead{\thepage}
\renewcommand{\headrulewidth}{0.4pt}

% ---------- 标题设置 ----------
\usepackage{titlesec}
\titleformat{\section}{\Large\bfseries}{}{0em}{}
\titleformat{\subsection}{\large\bfseries}{}{0em}{}
\titlespacing{\section}{0pt}{12pt}{6pt}
\titlespacing{\subsection}{0pt}{8pt}{4pt}

% ---------- 段落 ----------
\setlength{\parindent}{2em}
\setlength{\parskip}{0.5ex}

% ---------- 代码块 ----------
\usepackage{listings}
\usepackage{xcolor}
\lstset{
  language=Python,
  basicstyle=\small\ttfamily,
  numbers=left,
  numberstyle=\tiny,
  stepnumber=1,
  numbersep=5pt,
  backgroundcolor=\color{gray!10},
  frame=single,
  breaklines=true,
  captionpos=b,
  keywordstyle=\color{blue},
  commentstyle=\color{green!60!black},
  stringstyle=\color{red}
}

% ============================================
\begin{document}

% ---------- 首页：题目 + 摘要 + 关键词 ----------
\begin{center}
  {\LARGE\bfseries 山西大学第二十届数学建模竞赛}\\[0.8cm]
  {\huge\bfseries 题号：A（或B）}\\[0.5cm]
  {\LARGE\bfseries （论文题目）}\\[0.3cm]
  \rule{\textwidth}{0.5pt}\\[0.3cm]
\end{center}

\begin{center}
  {\large\bfseries 摘要}
\end{center}

\begin{quotation}
\noindent
%% ===== 摘要内容开始 =====
%% 要求：说明模型的主要特点、建模方法和主要结果
本文针对\underline{\hspace{4cm}}问题，综合运用\underline{\hspace{4cm}}
等方法，建立了\underline{\hspace{4cm}}模型。

针对问题一，\underline{\hspace{6cm}}。采用\underline{\hspace{3cm}}
方法求解，得到\underline{\hspace{3cm}}。

针对问题二，\underline{\hspace{6cm}}。采用\underline{\hspace{3cm}}
方法求解，得到\underline{\hspace{3cm}}。

\hspace{0pt}{\hfill（正文共\underline{\hspace{1cm}}页，附录\underline{\hspace{1cm}}页）}
%% ===== 摘要内容结束 =====
\end{quotation}

\vspace{0.3cm}
\noindent\textbf{关键词：} 关键词一；关键词二；关键词三；关键词四

\vspace{0.5cm}
\noindent\textbf{参赛队员：} 姓名1，姓名2，姓名3

\noindent\textbf{指导教师：} 教师姓名

\newpage

% ---------- 正文 ----------
\setcounter{page}{1}

\section{问题重述}

\subsection{问题背景}

% 用自己的话描述题目背景（不要照抄题目原文）

\subsection{问题要求}

% 列出题目要求解决的具体问题

\section{模型假设}

\begin{enumerate}
  \item \underline{\hspace{10cm}}
  \item \underline{\hspace{10cm}}
  \item \underline{\hspace{10cm}}
\end{enumerate}

\vspace{0.3cm}
\noindent{\small 注：每个假设后最好加一句理由说明。}

\section{符号说明}

\begin{center}
\begin{tabular}{ccc}
  \toprule
  \textbf{符号} & \textbf{含义} & \textbf{单位} \\
  \midrule
  $x_1$ & 变量一 &   \\
  $x_2$ & 变量二 &   \\
  $y$   & 因变量 &   \\
  $\alpha$ & 参数一 &   \\
  $\beta$  & 参数二 &   \\
  \bottomrule
\end{tabular}
\end{center}

\section{模型的建立与求解}

%% 每个子问题按：分析 -> 建模 -> 求解 -> 结果展示

\subsection{问题一的分析与模型建立}

\textbf{问题分析：}
问题一要解决\underline{\hspace{4cm}}，其核心在于\underline{\hspace{4cm}}。

\textbf{模型建立：}
基于上述分析，我们建立了如下模型：

\begin{equation}
  \label{eq:model1}
  \min Z = \sum_{i=1}^{n} c_i x_i
\end{equation}

\begin{align}
  \text{s.t.} \quad & \sum_{i=1}^{n} a_{ij} x_i \leq b_j,\quad j=1,\dots,m \\
  & x_i \geq 0,\quad i=1,\dots,n
\end{align}

其中，$c_i$表示\underline{\hspace{2cm}}，$a_{ij}$表示\underline{\hspace{2cm}}，
$b_j$表示\underline{\hspace{2cm}}。

\textbf{模型求解：}
使用MATLAB中的\texttt{linprog}函数求解，关键代码如下：

\begin{lstlisting}[caption=核心求解代码]
f = [c1, c2, ..., cn];
A = [a11, a12, ...; ...];
b = [b1, b2, ..., bm];
[x, fval] = linprog(f, A, b);
disp(['最优解: ', num2str(x')]);
disp(['最优值: ', num2str(fval)]);
\end{lstlisting}

\textbf{结果分析：}
求解结果如表~\ref{tab:result1}所示。

\begin{table}[H]
\centering
\caption{模型求解结果}
\label{tab:result1}
\begin{tabular}{cccc}
  \toprule
  参数 & 方案一 & 方案二 & 方案三 \\
  \midrule
  指标1 &  &  &  \\
  指标2 &  &  &  \\
  指标3 &  &  &  \\
  \bottomrule
\end{tabular}
\end{table}

\begin{figure}[H]
\centering
% \includegraphics[width=0.7\textwidth]{figure1.png}
\caption{结果可视化图}
\label{fig:result1}
\end{figure}

\newpage

\subsection{问题二的分析与模型建立}

\textbf{问题分析：}
问题二涉及\underline{\hspace{4cm}}，与问题一不同的是\underline{\hspace{4cm}}。

\textbf{模型建立：}
采用\underline{\hspace{3cm}}模型。

\begin{equation}
  \label{eq:model2}
  Y = \beta_0 + \beta_1 X_1 + \beta_2 X_2 + \cdots + \beta_k X_k + \varepsilon
\end{equation}

\section{结果分析与检验}

\subsection{灵敏度分析}

将关键参数$\alpha$在$[-20\%, +20\%]$范围内变化，观察结果变化。

\begin{figure}[H]
\centering
% \includegraphics[width=0.7\textwidth]{sensitivity.png}
\caption{灵敏度分析结果}
\label{fig:sensitivity}
\end{figure}

\subsection{误差分析}

\begin{table}[H]
\centering
\caption{误差分析}
\label{tab:error}
\begin{tabular}{cccc}
  \toprule
  样本 & 实际值 & 预测值 & 相对误差 \\
  \midrule
  1 &  &  & \\
  2 &  &  & \\
  3 &  &  & \\
  \bottomrule
\end{tabular}
\end{table}

\section{模型评价与改进}

\subsection{模型优点}
\begin{itemize}
  \item 模型结构清晰，易于理解和实现
  \item \underline{\hspace{8cm}}
\end{itemize}

\subsection{模型缺点}
\begin{itemize}
  \item \underline{\hspace{8cm}}
\end{itemize}

\subsection{改进方向}
\begin{itemize}
  \item 可以考虑引入\underline{\hspace{4cm}}来改进
\end{itemize}

\section{参考文献}

\begin{thebibliography}{9}
  \bibitem{ref1} 作者. 论文题目[J]. 期刊名, 出版年份, 卷(期): 起止页码.
  \bibitem{ref2} 作者. 书名[M]. 出版地: 出版社, 出版年份.
\end{thebibliography}

\newpage
\appendix
\section{附录：计算机程序}

%% 附件压缩包命名为：附件+姓名（如 附件_张三_李四_王五.zip）
%% 发送至 sxushuxuejianmo@163.com

\begin{lstlisting}[caption=主程序]
% 数学建模竞赛主程序
clear; clc; close all;
data = xlsread('data.xlsx');
% 模型求解
% 可视化
figure;
plot(x, y, 'b-o', 'LineWidth', 1.5);
xlabel('X');
ylabel('Y');
title('结果');
grid on;
\end{lstlisting}

\end{document}
